% Generated draft: requires mathematical and source review.
\documentclass[11pt]{article}
\usepackage[margin=25mm]{geometry}
\usepackage{fontspec}
\setmainfont{Latin Modern Roman}
\usepackage{amsmath,amssymb,mathtools,unicode-math}
\setmathfont{Latin Modern Math}
\usepackage{enumitem,needspace,xurl,hyperref}
\setlength{\parindent}{0pt}
\setlength{\parskip}{5pt}
\setlength{\emergencystretch}{3em}
\newcommand{\N}{\mathbb{N}}
\newcommand{\Z}{\mathbb{Z}}
\newcommand{\Q}{\mathbb{Q}}
\newcommand{\R}{\mathbb{R}}
\newcommand{\C}{\mathbb{C}}
\newcommand{\F}{\mathbb{F}}
\newcommand{\A}{\mathbb{A}}
\newcommand{\PP}{\mathbb{P}}
\newcommand{\E}{\mathbb{E}}
\newcommand{\eps}{\varepsilon}
\newcommand{\dd}{\,\mathrm{d}}
\begin{document}
\section*{UnsolvedMath ID: 700669}
\textbf{Source problem code:} catalogue entry
\subsection{Definitions and mathematical statement}
Regard a cardinal \(\mu\) as an index set, and let \(\Z^\mu\) be the abelian group of functions \(f\colon\mu\to\Z\), with pointwise addition. Define
\[
\operatorname{supp}(f)=\{\alpha\in\mu:f(\alpha)\ne0\},\qquad
\Z^{<\mu}=\{f\in\Z^\mu:|\operatorname{supp}(f)|<\mu\}.
\]
Whenever \(\Z^{<\mu}\) is a subgroup, define
\[
G_\mu=\Z^\mu/\Z^{<\mu}.
\]
In particular, this quotient is defined for every infinite cardinal \(\mu\). The source says only ``a cardinal number''; its unrestricted formulation requires qualification, since for finite \(\mu\ge2\) the displayed subset need not be closed under addition, and for \(\mu=0\) it is empty.

Let \(\omega_0\) denote the least infinite ordinal and \(\omega_1\) the least uncountable ordinal, used here as cardinal indices.
\begin{enumerate}
\item[(a)] Find a subgroup \(D\le G_{\omega_1}\) such that
\[
D\ne\{0\},\qquad D\not\cong G_{\omega_1},
\]
and \(D\) is a direct summand: there exists a subgroup \(E\le G_{\omega_1}\) satisfying
\[
G_{\omega_1}=D+E,\qquad D\cap E=\{0\}.
\]
\item[(b)] Investigate the structure of \(G_\mu\). The source supplies no particular classification criterion or more specific target for this investigation, and leaves the cardinal domain unqualified. For infinite cardinals the displayed construction is well-defined. The source additionally states that the structure of \(G_{\omega_0}\) is well known, without supplying a description or a separate reference for that assertion.
\end{enumerate}
\subsection{Short English statement}
Find a nonzero direct summand of \(G_{\omega_1}\) that is not isomorphic to the whole group, and investigate the groups \(G_\mu\) more generally. The source does not specify a precise classification target for the second part, and its cardinal domain needs clarification for finite cardinals.
\subsection{Sources}
\begin{itemize}
\item Supplied local UnsolvedMath registry record: permanent ID \(700669\), source problem code \texttt{catalogue entry}, title ``Kourovka Notebook Problem 10.54''; upstream revision \texttt{d371ee917ac9647fbb06cd6fc529deec6f3dda82}.
\item Supplied bibliographic source: \emph{The Kourovka Notebook (Unsolved Problems in Group Theory)}, original problem 10.54, 10th issue (1986); cited in the supplied record through the 21st issue edition (2026), \texttt{arXiv:1401.0300v48}. The supplied collection description names E. I. Khukhro and V. D. Mazurov as editors.
\item Attribution as supplied: the background and proposer field name S. V. Rychkov, with proposed year 1986; the statement itself begins with the attribution ``(R. G\"obel).'' The relationship between these attributions is not explained in the record.
\item The supplied registry labels the problem ``open'' and records its update as \texttt{2026-10-16T14:00:00Z}. This is supplied registry information, not an independently verified assessment of the literature or mathematical status. No independent source consultation was performed.
\end{itemize}
\end{document}
